\documentclass[11pt]{article}

\usepackage[margin=1in]{geometry}
\usepackage{amsmath,amssymb,amsthm,mathtools}
\usepackage{bm}
\usepackage{microtype}
\usepackage{hyperref}
\usepackage{listings}

\hypersetup{
  colorlinks=true,
  linkcolor=blue,
  citecolor=blue,
  urlcolor=blue
}

\lstdefinestyle{mystyle}{
  basicstyle=\ttfamily\small,
  breaklines=true,
  frame=single,
  columns=fullflexible
}
\lstset{style=mystyle}

\newtheorem{definition}{Definition}[section]
\newtheorem{theorem}[definition]{Theorem}
\newtheorem{proposition}[definition]{Proposition}
\newtheorem{lemma}[definition]{Lemma}
\newtheorem{corollary}[definition]{Corollary}
\newtheorem{remark}[definition]{Remark}

\newcommand{\R}{\mathbb{R}}
\newcommand{\N}{\mathbb{N}}
\newcommand{\norm}[1]{\left\lVert #1 \right\rVert}
\newcommand{\abs}[1]{\left\lvert #1 \right\rvert}

\title{Track C: WESAD CAS Proof-of-Concept \\
Defensive Prior Art Publication and Technical Report}

\author{Citizen Gardens \\ The Foundation of Multiplicity}

\date{April 27, 2026}

\begin{document}
\maketitle

\begin{abstract}

This document discloses a proof-of-concept application of the Codon-Aware Spectral (CAS) architecture to the WESAD wearable stress and affect detection dataset, constituting a Track~C extension of the Multiplicity Theory framework.\footnote{The WESAD dataset is a publicly available multimodal dataset for wearable stress and affect detection with 15 subjects and multiple physiological and motion channels, as described by Schmidt et al.\cite{Schmidt2018_WESAD,Schmidt2018_WESAD_UBI}}[web:25][web:408][web:405] The purpose of this publication is to establish prior art for the following contributions:

\begin{itemize}
  \item A CAS-style, channel-structured representation of WESAD signals using a fixed, low-dimensional axis system (Baseline, Stress, Amusement, and ancillary control channels).
  \item A Walsh--Hadamard-based spectral mixing layer on top of physiological features, treated as a CAS instance with orthogonal mixing and operator-norm guarantees.\cite{WHTOverview,WHTOrthogonal}[web:331][web:363][web:419]
  \item Integration of this spectral CAS representation with a Multiplicity-Crypto-compatible commitment and certification layer for stress/affect classification outputs (Track~B compatibility).
  \item A concrete Python reference implementation sketch demonstrating the full pipeline from raw WESAD signals to CAS features, classifier outputs, and commitment records, intended as a defense against subsequent patent claims over substantially similar architectures.\cite{WESAD_JohnsonGitHub,WESAD_MMIImplementation}[web:392][web:403]
\end{itemize}

The key novelty lies in treating WESAD as a CAS instance with (i) structured channels, (ii) orthogonal WHT-based spectral mixing, and (iii) optional attachment to multiplicity-crypto certification, thereby unifying stress-detection analytics with the broader Multiplicity Theory ecosystem.
\end{abstract}
\newpage
\tableofcontents
\newpage
\section{Background and Scope}

\subsection{WESAD Dataset}

The WESAD (Wearable Stress and Affect Detection) dataset comprises physiological and motion signals collected from 15 subjects via chest- and wrist-worn devices during lab sessions inducing neutral (baseline), stress, and amusement conditions.\cite{Schmidt2018_WESAD}[web:25][web:408][web:405] Available modalities include:

\begin{itemize}
  \item Electrocardiogram (ECG),
  \item Electrodermal activity (EDA),
  \item Electromyogram (EMG),
  \item Respiration (RESP),
  \item Skin temperature (TEMP),
  \item Blood volume pulse (BVP),
  \item Three-axis acceleration (ACC\_X, ACC\_Y, ACC\_Z).
\end{itemize}

The dataset has been used to benchmark stress vs.\ non-stress and multi-class affect detection (baseline vs.\ stress vs.\ amusement) using classical machine-learning and deep-learning models.\cite{Schmidt2018_WESAD,WESAD_JohnsonGitHub,WESAD_JaganGitHub,WESAD_Siyixu,WESAD_Srovins,WESAD_CrossModality,WESAD_CrossDomain}[web:408][web:399][web:393][web:395][web:396][web:22][web:418]

\subsection{CAS Architecture and WHT Layer}

In Track~A (Genomic CAS) and Track~B (Multiplicity-Crypto), CAS architectures used orthogonal Walsh--Hadamard-type transforms as mixing layers on top of structured channels.\cite{WHTOverview,WHTOrthogonal}[web:331][web:363][web:419] Track~C repurposes this machinery as follows:

\begin{enumerate}
  \item Define a small number of WESAD CAS channels capturing affective structure (baseline vs.\ stress vs.\ amusement), sensor-specific summary axes, and optional control axes (motion, artifacts).
  \item Map high-dimensional time-series features into these channels.
  \item Apply an orthogonal WHT mixing matrix to obtain a compact spectral representation.
  \item Use these CAS-spectral features as inputs to standard classifiers (e.g., logistic regression, SVM, shallow neural networks) for stress vs.\ non-stress and three-way affect detection.
  \item Optionally attach multiplicity-crypto commitments to model outputs and training events (Track~B integration).
\end{enumerate}

The remainder of this document states the CAS/WHT formulation for WESAD, states norm bounds, and provides a reference implementation sketch.

\section{Mathematical CAS Formulation for WESAD}

\subsection{Signal and Windowing Model}

Let \(s\) index subjects, \(c\) index sensor channels, and \(t\) denote discrete sample time.

\begin{definition}[Raw signal tensor]
Let
\[
  X_s(t) = \bigl( x_{s,c}(t) \bigr)_{c=1}^{C} \in \R^{C}
\]
be the multichannel time series for subject \(s\), where \(C\) aggregates all chosen WESAD modalities (e.g., ECG, EDA, RESP, TEMP, BVP, ACC components).\cite{Schmidt2018_WESAD,WESAD_CrossDomain}[web:408][web:418]
\end{definition}

We segment the signals into fixed-length windows.

\begin{definition}[Windowed segments]
Fix a window length \(L\) (in samples) and stride \(R\).
Define the \(j\)-th window for subject \(s\) as
\[
  X_{s,j} = \bigl(x_{s,c}(t)\bigr)_{c=1,\dots,C;\ t = t_j,\dots,t_j+L-1}
  \in \R^{C \times L},
\]
where \(t_j = t_0 + jR\).
Associate each window with a label \(y_{s,j} \in \{0,1,2\}\) for baseline, stress, and amusement, respectively.\cite{Schmidt2018_WESAD}[web:408][web:25]
\end{definition}

\subsection{Feature Extraction and Channel Aggregation}

For each sensor channel and window, we compute summary features such as mean, variance, spectral power, and domain-specific measures (e.g., HRV features from ECG, tonic/phasic components from EDA).\cite{WESAD_JohnsonGitHub,WESAD_MMIImplementation,WESAD_Siyixu}[web:392][web:403][web:395]

\begin{definition}[Feature vector]
For each window \((s,j)\), define a feature vector
\[
  f_{s,j} \in \R^d,
\]
built from concatenating features across all sensor modalities for that window.
\end{definition}

The CAS view is to group these features into a small number of channels.

\begin{definition}[Track C channel map]
Let \(\mathcal{K} = \{1,\dots,K\}\) index CAS channels.
Define disjoint index sets \(\mathcal{F}_k \subset \{1,\dots,d\}\) such that \(\bigcup_{k=1}^K \mathcal{F}_k = \{1,\dots,d\}\).
Define the channel aggregation map
\[
  A : \R^d \to \R^K, \quad
  (A f_{s,j})_k = \sum_{\ell \in \mathcal{F}_k} f_{s,j,\ell}.
\]
Then the channel vector for window \((s,j)\) is
\[
  a_{s,j} = A f_{s,j} \in \R^K.
\]
\end{definition}

Typical choices for channels include:

\begin{itemize}
  \item Cardio-respiratory channel (ECG, RESP-derived features),
  \item Electrodermal channel (EDA features),
  \item Motion channel (acceleration features),
  \item Temperature channel,
  \item Composite affective channels aligned with baseline, stress, and amusement axes (via supervised or unsupervised projection).
\end{itemize}

\subsection{Orthogonal WHT Mixing}

Let \(K = 2^m\) for some small \(m\), or let \(K\) be approximated by an orthogonal WHT-like matrix (as in Track~A) when not a power of two.\cite{WHTOverview,WHTOrthogonal}[web:331][web:363][web:419]

\begin{definition}[WESAD WHT mixing matrix]
Let \(H_K \in \R^{K \times K}\) be an orthogonal matrix (e.g., a normalized Walsh--Hadamard matrix or an orthogonal approximation when \(K\) is not a power of two).
Thus,
\[
  H_K H_K^\top = I_K = H_K^\top H_K.
\]
For each window \((s,j)\), define the CAS spectrum
\[
  h_{s,j} = H_K a_{s,j} \in \R^K.
\]
\end{definition}

\begin{proposition}[Norm preservation]
For all windows \((s,j)\),
\[
  \norm{h_{s,j}}_2 = \norm{a_{s,j}}_2.
\]
\end{proposition}

\begin{proof}
Because \(H_K\) is orthogonal, \(\norm{H_K x}_2 = \norm{x}_2\) for any \(x \in \R^K\).
This is standard for Walsh--Hadamard and other orthogonal transforms.\cite{WHTOverview,WHTOrthogonal}[web:331][web:363][web:419]
\end{proof}

\begin{proposition}[Aggregation operator norm bound]
If the aggregation matrix \(A\) satisfies \(\norm{A}_{2\to2} \le B_A\), then
\[
  \norm{h_{s,j}}_2 \le B_A \norm{f_{s,j}}_2.
\]
\end{proposition}

\begin{proof}
Since \(h_{s,j} = H_K A f_{s,j}\),
\[
  \norm{h_{s,j}}_2 \le \norm{H_K}_{2\to2} \norm{A}_{2\to2} \norm{f_{s,j}}_2 = B_A \norm{f_{s,j}}_2,
\]
because \(\norm{H_K}_{2\to2} = 1\) for an orthogonal matrix.
\end{proof}

Thus CAS-WHT features do not amplify norms beyond the aggregation operator norm, which is fixed given a channel design.

\section{Classification Layer and Certification Hooks}

\subsection{Feature Vectors for Stress/Affect Detection}

We consider one or more of the following feature families for stress/affect classification:

\begin{itemize}
  \item Raw aggregated channels \(a_{s,j}\),
  \item Spectral channels \(h_{s,j}\),
  \item Residuals or deviations from baseline (within-subject normalization),
  \item Concatenations of channel and spectral features.
\end{itemize}

\begin{definition}[CAS feature map]
Let \(\Phi : \R^d \to \R^D\) be the composition
\[
  \Phi(f_{s,j}) = \Psi\bigl( a_{s,j}, h_{s,j}\bigr),
\]
where \(\Psi\) is any fixed, differentiable map (e.g., stacking, residual computation, normalization).
The CAS feature for window \((s,j)\) is
\[
  z_{s,j} = \Phi(f_{s,j}) \in \R^D.
\]
\end{definition}

Because \(A\) and \(H_K\) are linear and bounded in operator norm, and \(\Psi\) can be designed to be Lipschitz, the overall map \(\Phi\) is Lipschitz, which simplifies stability and optimization analysis.

\subsection{Logistic Regression Model}

For binary stress vs.\ non-stress classification, we define:

\begin{definition}[Binary Track C classifier]
Let \(y_{s,j} \in \{0,1\}\) indicate non-stress vs.\ stress.
A logistic regression model uses parameters \(\theta \in \R^D\) and
\[
  p_{\theta}(y_{s,j}=1 \mid z_{s,j}) = \sigma(\theta^\top z_{s,j}),
\]
where \(\sigma(z) = 1/(1+e^{-z})\).
\end{definition}

The per-sample loss is
\[
  \ell(\theta; z_{s,j}, y_{s,j})
  = -y_{s,j} \log \sigma(\theta^\top z_{s,j})
    - (1-y_{s,j}) \log \bigl(1-\sigma(\theta^\top z_{s,j})\bigr).
\]

\begin{proposition}[Gradient bound]
If \(\norm{z_{s,j}}_2 \le B_Z\) for all \((s,j)\), then
\[
  \norm{\nabla_\theta \ell(\theta; z_{s,j}, y_{s,j})}_2 \le B_Z.
\]
\end{proposition}

\begin{proof}
The gradient is
\[
  \nabla_\theta \ell = \bigl(\sigma(\theta^\top z_{s,j}) - y_{s,j}\bigr) z_{s,j}.
\]
Because \(\sigma(\theta^\top z_{s,j}) \in (0,1)\) and \(y_{s,j} \in \{0,1\}\), we have \(\abs{\sigma(\theta^\top z_{s,j}) - y_{s,j}} \le 1\).
Thus
\[
  \norm{\nabla_\theta \ell}_2 \le \norm{z_{s,j}}_2 \le B_Z.
\]
\end{proof}

Using the previous norm bounds for \(A\) and \(H_K\), one can guarantee a fixed \(B_Z\) given bounded feature vectors \(f_{s,j}\).

\subsection{Certification Hooks (Track B Compatibility)}

To integrate with Track~B's multiplicity-crypto layer, Track~C can define:

\begin{itemize}
  \item A per-subject or per-session multiplicity operator \(M_t^{(C)}\) counting events such as model training epochs, threshold crossings, and classification outputs.
  \item Commitment values \(C(v,\rho)\) over CAS feature aggregates or model snapshots using Pedersen commitments in a BN254-like group.\cite{Pedersen_ZKDocs,Pedersen_RareSkills}[web:366][web:369]
  \item HKDF-based key derivation for encrypted logging of classification results, with WESAD CAS metadata included as associated data.
\end{itemize}

The detailed multiplicity-crypto machinery is beyond the scope of this appendix, but is structurally inherited from Track~B.

\section{Reference Implementation Sketch}

This section provides a Python-style code snippet for a minimal Track~C CAS pipeline on WESAD. Implementations in existing repositories demonstrate traditional pipelines, but do not use a CAS/WHT structure.\cite{WESAD_JohnsonGitHub,WESAD_MMIImplementation}[web:392][web:403]

\subsection{Data Loading and Preprocessing}

\begin{lstlisting}[language=Python,caption={Minimal WESAD CAS feature pipeline (conceptual).}]
import numpy as np
from scipy.signal import welch
from scipy.linalg import hadamard

# --- Configuration ---

WINDOW_SECONDS = 60
STRIDE_SECONDS = 30
FS_ECG = 700.0  # example sampling rate, see WESAD documentation
K_CHANNELS = 8  # number of CAS channels (power of 2 for standard WHT)

# --- Helper: windowing ---

def sliding_windows(x, window, stride):
    """Simple sliding-window generator along axis=0."""
    start = 0
    n = x.shape[0]
    while start + window <= n:
        yield x[start:start + window]
        start += stride

# --- Feature extraction per window ---

def extract_features_window(ecg, eda, resp, temp, acc):
    """Return feature vector f in R^d for one window."""
    feats = []

    # Example ECG features (time-domain + frequency-domain)
    feats.append(np.mean(ecg))
    feats.append(np.std(ecg))
    f_ecg, p_ecg = welch(ecg, fs=FS_ECG, nperseg=len(ecg))
    feats.append(np.trapz(p_ecg[(f_ecg >= 0.04) & (f_ecg <= 0.15)]) )  # LF power
    feats.append(np.trapz(p_ecg[(f_ecg >= 0.15) & (f_ecg <= 0.4)]) )   # HF power

    # EDA features
    feats.append(np.mean(eda))
    feats.append(np.std(eda))

    # Respiration features
    feats.append(np.mean(resp))
    feats.append(np.std(resp))

    # Temperature features
    feats.append(np.mean(temp))

    # ACC features (3-axis)
    for axis in range(acc.shape[1]):
        feats.append(np.mean(acc[:, axis]))
        feats.append(np.std(acc[:, axis]))

    return np.array(feats, dtype=np.float32)  # f in R^d

# --- Channel aggregation ---

def build_channel_matrix(d, K):
    """
    Build aggregation matrix A in R^{K x d}.
    Here we use a simple block assignment; in practice, this
    can be hand-crafted to group semantically related features.
    """
    A = np.zeros((K, d), dtype=np.float32)
    block = int(np.ceil(d / K))
    idx = 0
    for k in range(K):
        for j in range(block):
            if idx < d:
                A[k, idx] = 1.0
                idx += 1
    return A

# --- WHT mixing ---

def hadamard_mixing(K):
    """
    Return orthogonal mixing matrix H_K in R^{K x K}.
    For K a power of 2, use normalized Hadamard matrix.
    """
    H = hadamard(K).astype(np.float32)
    H /= np.sqrt(K)
    return H

# --- Main CAS feature pipeline for one subject ---

def cas_features_for_subject(subj_data):
    """
    subj_data: dict with keys 'ecg', 'eda', 'resp', 'temp', 'acc', 'labels'
      Each signal is a 1D or 2D numpy array aligned in time.
    Returns:
      Z: CAS features in R^{num_windows x D}
      Y: integer labels per window
    """
    ecg = subj_data["ecg"]
    eda = subj_data["eda"]
    resp = subj_data["resp"]
    temp = subj_data["temp"]
    acc = subj_data["acc"]  # shape: (T, 3)
    labels = subj_data["labels"]  # per-sample label

    win = int(WINDOW_SECONDS * FS_ECG)
    stride = int(STRIDE_SECONDS * FS_ECG)

    # Build A and H_K once
    # First, get d from a dummy window
    f0 = extract_features_window(ecg[:win], eda[:win], resp[:win], temp[:win], acc[:win])
    d = f0.shape[0]
    A = build_channel_matrix(d, K_CHANNELS)
    H = hadamard_mixing(K_CHANNELS)

    Z = []
    Y = []

    for w_idx, (ecg_w, eda_w, resp_w, temp_w, acc_w) in enumerate(
        zip(sliding_windows(ecg, win, stride),
            sliding_windows(eda, win, stride),
            sliding_windows(resp, win, stride),
            sliding_windows(temp, win, stride),
            sliding_windows(acc, win, stride))
    ):
        f = extract_features_window(ecg_w, eda_w, resp_w, temp_w, acc_w)   # in R^d
        a = A @ f                                                           # in R^K
        h = H @ a                                                           # in R^K

        # Simple CAS feature: concatenate a and h
        z = np.concatenate([a, h], axis=0)
        Z.append(z)

        # Majority label within window
        label_w = np.bincount(labels[w_idx * stride: w_idx * stride + win]).argmax()
        Y.append(label_w)

    return np.stack(Z, axis=0), np.array(Y, dtype=np.int64)
\end{lstlisting}

This code demonstrates:

\begin{itemize}
  \item Construction of windowed WESAD features from ECG, EDA, RESP, TEMP, and ACC.\cite{WESAD_JohnsonGitHub,WESAD_MMIImplementation}[web:392][web:403]
  \item A channel aggregation matrix \(A\) and orthogonal WHT mixing \(H\) as in the CAS formalism.
  \item A final CAS feature vector per window \(z_{s,j} \in \R^{2K}\).
\end{itemize}

\subsection{Classifier and Certification Hook}

\begin{lstlisting}[language=Python,caption={Classifier training and placeholder for certification.}]
from sklearn.linear_model import LogisticRegression
from sklearn.model_selection import train_test_split
from sklearn.metrics import classification_report

# Suppose we aggregate all subjects:
Z_all, Y_all = [], []
for subj_data in all_subjects_data:  # list of subject dicts
    Z_subj, Y_subj = cas_features_for_subject(subj_data)
    Z_all.append(Z_subj)
    Y_all.append(Y_subj)
Z_all = np.concatenate(Z_all, axis=0)
Y_all = np.concatenate(Y_all, axis=0)

# Binary label example: stress vs. non-stress
Y_bin = (Y_all == STRESS_LABEL).astype(int)

X_train, X_test, y_train, y_test = train_test_split(
    Z_all, Y_bin, test_size=0.2, stratify=Y_bin, random_state=42
)

clf = LogisticRegression(max_iter=1000)
clf.fit(X_train, y_train)

y_pred = clf.predict(X_test)
print(classification_report(y_test, y_pred))

# --- Certification hook (conceptual) ---

def pedersen_commit(value_bytes, blinding_scalar):
    # placeholder stub for BN254 Pedersen commitment
    # using external library in a real implementation
    # C = rho * G + v * H
    pass

# Example: commit to model parameters
theta_bytes = clf.coef_.astype(np.float32).tobytes() + clf.intercept_.astype(np.float32).tobytes()
rho = os.urandom(32)  # real implementation: map to scalar mod r
model_commitment = pedersen_commit(theta_bytes, rho)
\end{lstlisting}

In a full Track~C system, the model commitment would be bound to training transcript and evaluation results using the Track~B multiplicity-crypto scheme (HKDF-derived keys, multiplicity operator, and transcript chains).

\section{Prior Art Scope and Claims}

This document intends to place the following objects and constructions into the public domain as prior art:

\begin{enumerate}
  \item \textbf{CAS representation for WESAD}: The notion of defining a finite set of WESAD CAS channels, aggregating high-dimensional time-series features into these channels, and applying an orthogonal Walsh--Hadamard or WHT-like transform to obtain a low-dimensional spectral representation for stress/affect detection.\cite{Schmidt2018_WESAD,WHTOverview,WHTOrthogonal}[web:408][web:331][web:363][web:419]
  \item \textbf{Operator-norm guarantees}: The use of explicit operator-norm bounds on the channel aggregation matrix and orthogonal transform to guarantee stability and boundedness of CAS features when used in logistic or related classifiers.
  \item \textbf{Track~B-compatible certification hooks}: The explicit inclusion of multiplicity-counting operators and Pedersen-style commitments on top of WESAD CAS features and model parameters, as a pattern for certified stress/affect analytics, even when not fully implemented in this proof-of-concept.\cite{Pedersen_ZKDocs,Pedersen_RareSkills}[web:366][web:369]
  \item \textbf{Reference implementation}: The presented Python pipeline structure from raw WESAD data to CAS features, classifier training, and commitment hooks, as a specific algorithmic template.
\end{enumerate}

These items, together with the mathematical derivations and code sketches given above, should be considered prior art against any attempt to claim exclusive rights over substantially similar combinations of WESAD, CAS, WHT, and multiplicity-crypto certification.

\appendix

\section{Signal Model and Windowed Feature Space}

\subsection{WESAD Multichannel Signals}

Track~C operates on multichannel time-series signals derived from the WESAD dataset (ECG, EDA, RESP, TEMP, ACC, etc.), but the mathematics below assumes a general multichannel sequence model.

\begin{definition}[Raw multichannel signal]
Let \(C \in \N\) be the number of sensor channels.
For a given subject \(s\), define the discrete-time multichannel signal
\[
  X_s(t) = \bigl(x_{s,c}(t)\bigr)_{c=1}^C \in \R^C, \quad t = 0,1,2,\dots .
\]
\end{definition}

\begin{definition}[Windowed segments]
Fix a window length \(L \in \N\) and stride \(R \in \N\).
The \(j\)-th window for subject \(s\) is the tensor
\[
  X_{s,j} = \bigl(x_{s,c}(t)\bigr)_{c=1,\dots,C;\, t = t_j,\dots,t_j+L-1}
  \in \R^{C \times L},
\]
where \(t_j = t_0 + jR\).
\end{definition}

Each window is associated with an affect label \(y_{s,j}\) (e.g.~baseline, stress, amusement), but the label is not needed for the norm bounds.

\subsection{Feature Extraction}

\begin{definition}[Per-window feature vector]
For each window \((s,j)\), a deterministic feature extractor
\[
  \mathcal{F} : \R^{C \times L} \to \R^d
\]
produces a feature vector
\[
  f_{s,j} = \mathcal{F}(X_{s,j}) \in \R^d.
\]
\end{definition}

\begin{remark}
In practice, \(\mathcal{F}\) concatenates statistics such as means, variances, frequency-band powers, and domain-specific features (e.g.~HRV from ECG, tonic/phasic components from EDA). The derivations below treat \(\mathcal{F}\) as an arbitrary, fixed map and place assumptions only on the resulting \(f_{s,j}\).
\end{remark}

\section{Channel Aggregation and Operator Norms}

\subsection{Channel Aggregation Matrix}

The CAS view groups the \(d\) features into \(K\) channels.

\begin{definition}[Aggregation matrix]
Let \(K \in \N\) be the number of CAS channels.
The channel aggregation matrix is
\[
  A \in \R^{K \times d},
\]
and the channel-intensity vector for window \((s,j)\) is
\[
  a_{s,j} = A f_{s,j} \in \R^K.
\]
\end{definition}

No special structure is required, but a common pattern is that each row of \(A\) selects and sums features corresponding to a particular modality or semantic grouping.

\begin{definition}[Operator norms]
For a matrix \(M \in \R^{m \times n}\),
\[
  \norm{M}_{2\to 2} = \sup_{x \ne 0} \frac{\norm{M x}_2}{\norm{x}_2}
\]
denotes the spectral (operator 2-) norm. We also use the Frobenius norm
\[
  \norm{M}_F = \biggl( \sum_{i,j} M_{ij}^2 \biggr)^{1/2}
\]
when convenient.
\end{definition}

\begin{proposition}[Channel norm bound]
Suppose \(\norm{A}_{2\to 2} \le B_A\) for some \(B_A > 0\).
Then for any feature vector \(f \in \R^d\),
\[
  \norm{A f}_2 \le B_A \norm{f}_2.
\]
In particular, for each window \((s,j)\),
\[
  \norm{a_{s,j}}_2 \le B_A \norm{f_{s,j}}_2.
\]
\end{proposition}

\begin{proof}
By the definition of the operator 2-norm,
\[
  \norm{A f}_2 \le \norm{A}_{2\to2} \norm{f}_2 \le B_A \norm{f}_2.
\]
Substituting \(f = f_{s,j}\) yields the second statement.
\end{proof}

\subsection{Specific Structures of \texorpdfstring{$A$}{A} and Norm Estimates}

If each feature participates in at most one channel and all nonzero entries of \(A\) are equal to 1, then we can bound \(\norm{A}_{2\to2}\) in terms of row sparsity.

\begin{proposition}[Row-sparsity bound]
Assume \(A \in \{0,1\}^{K \times d}\) and that each column has at most one nonzero entry.
Let
\[
  r_k = \#\{ \ell : A_{k,\ell} = 1 \}, \quad k=1,\dots,K,
\]
and let \(r_{\max} = \max_k r_k\).
Then
\[
  \norm{A}_{2\to2} \le \sqrt{r_{\max}}.
\]
\end{proposition}

\begin{proof}
Compute
\[
  (A A^\top)_{k,k'} = \sum_{\ell=1}^d A_{k,\ell} A_{k',\ell}.
\]
Because each column has at most one nonzero, \(A_{k,\ell} A_{k',\ell} = 0\) whenever \(k \ne k'\).
Thus \(A A^\top\) is diagonal with
\[
  (A A^\top)_{k,k} = \sum_{\ell=1}^d A_{k,\ell}^2 = r_k.
\]
The eigenvalues of \(A A^\top\) are therefore \(\{r_k\}_{k=1}^K\).
The spectral norm of \(A\) satisfies
\[
  \norm{A}_{2\to2}^2 = \lambda_{\max}(A A^\top) = r_{\max},
\]
so \(\norm{A}_{2\to2} \le \sqrt{r_{\max}}\).
\end{proof}

\begin{corollary}[Norm bound with row-sparsity]
If \(A\) satisfies the assumptions above, then for each window \((s,j)\),
\[
  \norm{a_{s,j}}_2 \le \sqrt{r_{\max}} \,\norm{f_{s,j}}_2.
\]
\end{corollary}

\section{Orthogonal WHT Mixing Layer}

\subsection{Walsh--Hadamard-Type Matrix}

Track~C uses an orthogonal mixing matrix \(H_K\), often chosen as a normalized Walsh--Hadamard matrix when \(K\) is a power of two.

\begin{definition}[Orthogonal mixing matrix]
Let \(H_K \in \R^{K \times K}\) be a real orthogonal matrix:
\[
  H_K H_K^\top = I_K = H_K^\top H_K.
\]
Given a channel vector \(a_{s,j}\), the mixed spectral vector is
\[
  h_{s,j} = H_K a_{s,j} \in \R^K.
\]
\end{definition}

\begin{proposition}[Norm preservation under \(H_K\)]
For any \(x \in \R^K\),
\[
  \norm{H_K x}_2 = \norm{x}_2.
\]
In particular, for each window \((s,j)\),
\[
  \norm{h_{s,j}}_2 = \norm{a_{s,j}}_2.
\]
\end{proposition}

\begin{proof}
Orthogonality implies
\[
  \norm{H_K x}_2^2 = (H_K x)^\top (H_K x)
  = x^\top H_K^\top H_K x
  = x^\top x = \norm{x}_2^2.
\]
Taking square roots yields \(\norm{H_K x}_2 = \norm{x}_2\).
Substituting \(x = a_{s,j}\) gives the claim about \(\norm{h_{s,j}}_2\).
\end{proof}

\subsection{Composite Linear Map and Operator Norms}

The overall linear map from feature space to spectral space is
\[
  M = H_K A : \R^d \to \R^K, \quad M f = H_K (A f).
\]

\begin{proposition}[Operator norm of \(M = H_K A\)]
The operator 2-norm of \(M\) equals that of \(A\):
\[
  \norm{M}_{2\to2} = \norm{A}_{2\to2}.
\]
\end{proposition}

\begin{proof}
Because \(H_K\) is orthogonal, \(\norm{H_K x}_2 = \norm{x}_2\) for all \(x\).
Thus for any \(f\),
\[
  \norm{M f}_2 = \norm{H_K A f}_2 = \norm{A f}_2.
\]
Taking suprema over nonzero \(f\) yields
\[
  \norm{M}_{2\to2} = \sup_{f \ne 0} \frac{\norm{M f}_2}{\norm{f}_2}
  = \sup_{f \ne 0} \frac{\norm{A f}_2}{\norm{f}_2}
  = \norm{A}_{2\to2}.
\]
\end{proof}

\begin{corollary}[Spectral norm bound for Track C]
If \(\norm{A}_{2\to2} \le B_A\), then for each window \((s,j)\),
\[
  \norm{h_{s,j}}_2 \le B_A \norm{f_{s,j}}_2.
\]
\end{corollary}

\begin{proof}
We have \(h_{s,j} = M f_{s,j}\).
Thus
\[
  \norm{h_{s,j}}_2 = \norm{M f_{s,j}}_2 \le \norm{M}_{2\to2} \norm{f_{s,j}}_2 = \norm{A}_{2\to2} \norm{f_{s,j}}_2 \le B_A \norm{f_{s,j}}_2.
\]
\end{proof}

\section{CAS Feature Vector and Boundedness}

\subsection{Definition of Track C CAS Feature}

For each window, Track~C defines a CAS feature vector as some function of \(a_{s,j}\) and \(h_{s,j}\).

\begin{definition}[CAS feature map]
Let \(\Psi : \R^K \times \R^K \to \R^D\) be a fixed (possibly nonlinear) map.
The Track~C CAS feature for window \((s,j)\) is
\[
  z_{s,j} = \Phi(f_{s,j}) := \Psi\bigl(a_{s,j}, h_{s,j}\bigr)
  = \Psi\bigl(A f_{s,j}, H_K A f_{s,j}\bigr) \in \R^D.
\]
\end{definition}

Examples include simple concatenation \(\Psi(a,h) = (a,h)\), residuals, or normalized combinations.

\begin{definition}[Lipschitz constant of \(\Psi\)]
Suppose there exists \(L_\Psi > 0\) such that for all \((a,h)\) and \((a',h')\),
\[
  \norm{\Psi(a,h) - \Psi(a',h')}_2
  \le L_\Psi \bigl( \norm{a-a'}_2 + \norm{h-h'}_2 \bigr).
\]
\end{definition}

\begin{proposition}[Lipschitz constant of \(\Phi\)]
Assume \(\norm{A}_{2\to2} \le B_A\) and that \(\Psi\) is Lipschitz with constant \(L_\Psi\) as above.
Then the composite map \(\Phi : \R^d \to \R^D\) satisfies
\[
  \norm{\Phi(f) - \Phi(f')}_2 \le 2 B_A L_\Psi \norm{f - f'}_2
\]
for all \(f,f' \in \R^d\).
\end{proposition}

\begin{proof}
Write \(a = A f\), \(h = H_K A f\), and similarly for \(a',h'\).
Since \(H_K\) is orthogonal,
\[
  \norm{h - h'}_2 = \norm{H_K(A f - A f')}_2 = \norm{A(f-f')}_2 \le B_A \norm{f-f'}_2,
\]
and also \(\norm{a - a'}_2 = \norm{A(f-f')}_2 \le B_A \norm{f-f'}_2\).
Therefore,
\begin{align*}
  \norm{\Phi(f) - \Phi(f')}_2
  &= \norm{\Psi(a,h) - \Psi(a',h')}_2 \\
  &\le L_\Psi \bigl( \norm{a-a'}_2 + \norm{h-h'}_2 \bigr) \\
  &\le L_\Psi (B_A \norm{f-f'}_2 + B_A \norm{f-f'}_2) \\
  &= 2 B_A L_\Psi \norm{f-f'}_2.
\end{align*}
\end{proof}

\begin{corollary}[Uniform bound given bounded feature norm]
If \(\norm{f_{s,j}}_2 \le B_f\) for all \((s,j)\), then
\[
  \norm{a_{s,j}}_2 \le B_A B_f, \quad
  \norm{h_{s,j}}_2 \le B_A B_f,
\]
and, in particular, if \(\Psi\) is Lipschitz and \(\Psi(0,0) = 0\),
\[
  \norm{z_{s,j}}_2 \le 2 B_A L_\Psi B_f.
\]
\end{corollary}

\section{Logistic Regression on CAS Features}

\subsection{Model Definition}

For binary stress vs.\ non-stress, each window \((s,j)\) is labeled \(y_{s,j} \in \{0,1\}\).

\begin{definition}[Logistic regression on CAS features]
Let \(\theta \in \R^D\) be model parameters.
Define
\[
  p_\theta(y_{s,j}=1 \mid z_{s,j}) = \sigma(\theta^\top z_{s,j}),
\]
where \(\sigma(z) = \frac{1}{1+e^{-z}}\).
The per-sample negative log-likelihood is
\[
  \ell(\theta; z_{s,j}, y_{s,j})
  = -y_{s,j} \log \sigma(\theta^\top z_{s,j})
    - (1-y_{s,j}) \log \bigl(1-\sigma(\theta^\top z_{s,j})\bigr).
\]
\end{definition}

\subsection{Gradient Norm Bound}

\begin{proposition}[Gradient bound]
If \(\norm{z_{s,j}}_2 \le B_Z\) for all \((s,j)\), then for all \(\theta\),
\[
  \norm{\nabla_\theta \ell(\theta; z_{s,j}, y_{s,j})}_2 \le B_Z.
\]
\end{proposition}

\begin{proof}
The gradient is
\[
  \nabla_\theta \ell(\theta; z_{s,j}, y_{s,j})
  = \bigl(\sigma(\theta^\top z_{s,j}) - y_{s,j}\bigr) z_{s,j}.
\]
Since \(y_{s,j} \in \{0,1\}\) and \(\sigma(\theta^\top z_{s,j}) \in (0,1)\), we have
\(\abs{\sigma(\theta^\top z_{s,j}) - y_{s,j}} \le 1\).
Therefore,
\[
  \norm{\nabla_\theta \ell}_2
  \le \abs{\sigma(\theta^\top z_{s,j}) - y_{s,j}} \,\norm{z_{s,j}}_2
  \le \norm{z_{s,j}}_2 \le B_Z.
\]
\end{proof}

\begin{corollary}[Lipschitz gradient]
If \(B_Z < \infty\), the logistic loss has gradients uniformly bounded in norm by \(B_Z\), which facilitates step-size selection and convergence guarantees for first-order optimization methods.
\end{corollary}

\section{Within-Subject Normalization and Baseline Deviations}

Track~C can also use within-subject baseline normalization for CAS features.

\subsection{Subject Baseline CAS Feature}

\begin{definition}[Subject baseline CAS feature]
For subject \(s\), let \(\mathcal{B}_s\) denote the set of baseline windows.
Define the baseline CAS mean as
\[
  \bar{z}_s^{(\mathrm{base})}
  = \frac{1}{\abs{\mathcal{B}_s}} \sum_{j \in \mathcal{B}_s} z_{s,j}.
\]
\end{definition}

\begin{definition}[Deviation from baseline]
For any window \((s,j)\), define the deviation
\[
  \Delta z_{s,j} = z_{s,j} - \bar{z}_s^{(\mathrm{base})}.
\]
\end{definition}

\subsection{Norm Bounds for Deviations}

\begin{proposition}[Deviation norm bound]
If \(\norm{z_{s,j}}_2 \le B_Z\) for all windows of subject \(s\), then
\[
  \norm{\Delta z_{s,j}}_2 \le 2 B_Z
\]
for all \(j\).
\end{proposition}

\begin{proof}
Using the triangle inequality,
\[
  \norm{\Delta z_{s,j}}_2
  = \norm{z_{s,j} - \bar{z}_s^{(\mathrm{base})}}_2
  \le \norm{z_{s,j}}_2 + \norm{\bar{z}_s^{(\mathrm{base})}}_2.
\]
Moreover,
\[
  \norm{\bar{z}_s^{(\mathrm{base})}}_2
  = \biggl\lVert \frac{1}{\abs{\mathcal{B}_s}} \sum_{j \in \mathcal{B}_s} z_{s,j} \biggr\rVert_2
  \le \frac{1}{\abs{\mathcal{B}_s}} \sum_{j \in \mathcal{B}_s} \norm{z_{s,j}}_2
  \le B_Z.
\]
Hence \(\norm{\Delta z_{s,j}}_2 \le B_Z + B_Z = 2 B_Z\).
\end{proof}

\section{Summary of Operator Norm Bounds}

For the Track~C WESAD CAS proof-of-concept, the key bounds are:

\begin{itemize}
  \item \textbf{Aggregation matrix:}
    \[
      \norm{A}_{2\to2} \le B_A, \quad
      \norm{a_{s,j}}_2 \le B_A \norm{f_{s,j}}_2.
    \]
    If \(A\) is a 0--1 matrix with each column used at most once and row sparsity \(r_{\max}\), then \(\norm{A}_{2\to2} \le \sqrt{r_{\max}}\).
  \item \textbf{Orthogonal mixing:}
    \[
      \norm{H_K x}_2 = \norm{x}_2, \quad
      \norm{h_{s,j}}_2 = \norm{a_{s,j}}_2.
    \]
  \item \textbf{Composite map:}
    \[
      M = H_K A, \quad
      \norm{M}_{2\to2} = \norm{A}_{2\to2}.
    \]
  \item \textbf{CAS feature map:}
    If \(\Psi\) is Lipschitz with constant \(L_\Psi\) and \(\norm{A}_{2\to2} \le B_A\), then
    \[
      \norm{\Phi(f) - \Phi(f')}_2 \le 2 B_A L_\Psi \norm{f-f'}_2.
    \]
    With bounded feature norm \(\norm{f_{s,j}}_2 \le B_f\), all CAS features satisfy \(\norm{z_{s,j}}_2 \le 2 B_A L_\Psi B_f\).
  \item \textbf{Logistic regression gradient:}
    If \(\norm{z_{s,j}}_2 \le B_Z\), then
    \[
      \norm{\nabla_\theta \ell(\theta; z_{s,j}, y_{s,j})}_2 \le B_Z
    \]
    for all \(\theta\) and all windows.
  \item \textbf{Baseline deviations:}
    With \(\norm{z_{s,j}}_2 \le B_Z\), the within-subject deviations satisfy
    \[
      \norm{\Delta z_{s,j}}_2 \le 2 B_Z.
    \]
\end{itemize}

These results collectively provide the explicit proofs and norm bounds underpinning Track~C’s WESAD CAS formulation, ensuring that the spectral and classifier layers are numerically stable and analytically well-controlled.


\nocite{*}
\bibliographystyle{unsrt}  
\bibliography{references}  


\end{document}
